% HD-S34. GPT-6 Astra Pro; inherited effort not exposed; /root/sixfolds.
% 21 November 2026. Actual K-defined one-dimensional blocks only.
\subsection{HD-S34: every natural Gaussian factor with \(K\)-defined monomial lines is split}
\label{sec:gaussian-monomial-scope}

Let \(K=\mathbb Q(i)\), and let \(\rho:G\to GL_K(V)\)
be an absolutely irreducible finite representation whose character
field is exactly \(K\). Suppose it has an actual \(K\)-defined
system of one-dimensional blocks
\[
                    V=\bigoplus_{j=1}^m L_j,
\]
permuted by \(G\). For every connected sphere cover with full
monodromy \(G\), the natural primitive Jacobian factor of
signature \((3,3)\), with its restricted theta polarization,
has ordinary Hermitian determinant class \([-1]\), and hence
is split by the explicit isometry construction of HD-S03a.
There is no restriction on its branch profile or period rank.
The block system here is defined over \(K\); a decomposition
available only after extending the coefficient field is not
included in the statement.

\paragraph{1. An actual coset basis with fourth-root entries.}
Irreducibility makes the action on the block set transitive,
since the sum over any orbit is an invariant subspace. Put
\(H=\operatorname{Stab}_G(L_1)\), choose \(0\ne v_1\in L_1\),
and choose elements \(t_j\) taking \(L_1\) to \(L_j\), with
\(t_1=1\). Set \(v_j=\rho(t_j)v_1\). These vectors form a
\(K\)-basis. Let \(\chi:H\to K^\times\) be the actual
one-dimensional action on \(L_1\). Every \(\chi(h)\) is a
root of unity, since \(h\) has finite order. A root of unity in
\(K\) is an integral Gaussian number of norm one, so it is one
of \(1,-1,i,-i\).

For \(g\in G\) and a block index \(j\), the element
\(h_{g,j}=t_{g(j)}^{-1}gt_j\) lies in \(H\), and the actual
matrix formula is
\[
             \rho(g)v_j=\chi(h_{g,j})v_{g(j)}.
 \tag{S34.1}
\]
Thus the matrix whose columns are the \(v_j\) gives the precise
\(K\)-linear change to a unit monomial representation. In this
basis \(h_0(v_j,v_k)=\delta_{jk}\) is a positive invariant
Hermitian form, and
\[
                 L=\bigoplus_j Rv_j,\qquad
                 R=\mathbb Z[i,1/2]
 \tag{S34.2}
\]
is an invariant self-dual lattice.

The character \(\chi\) has image \(\mu_4\). Indeed, if its
image were contained in \(\{1,-1\}\), every matrix in (S34.1)
would have rational trace, contrary to the stated character field.
Choose and fix \(h_0'\in H\) with \(\chi(h_0')=i\).
This is a block-stabilizer element; it need not be central in
\(G\).

For any other specified positive invariant \(K\)-Hermitian form,
the same basis change carries it to \(aI_m\), where
\(a\in\mathbb Q_{>0}\). To check this, its Gram matrix relative
to \(I_m\) commutes with every monomial matrix; absolute
irreducibility makes it a scalar, and Hermitianity puts that scalar
in \(K\cap\mathbb R\). Thus comparison with such a coefficient
form is through an explicit basis change and positive rational
scalar. Its rank-six determinant effect is a norm.

\paragraph{2. Local cycle modules and global coinvariants.}
For a local monodromy \(g\), take one cycle of its permutation
of the lines, of length \(\ell\), starting at \(L_a\). The
vectors \(w_j=\rho(g)^jv_a\), \(0\le j<\ell\), are unit
multiples of the original cycle basis. In them
\[
 g w_j=w_{j+1}\quad(j<\ell-1),\qquad
 g w_{\ell-1}=\zeta w_0,\qquad \zeta\in\mu_4.
\]
The invariant and coinvariant calculation is exactly the one
proved in HD-S33: the coinvariants are \(R/(\zeta-1)\).
For \(\zeta=1\), both modules are free of rank one, with
invariant generator \(\sum w_j\), coinvariant generator
\([w_0]\), and pairing one. For \(\zeta\ne1\), both vanish,
since \(\zeta-1\) is a unit of \(R\); the geometric-sum
matrix is an explicit inverse to \(1-g\). Thus every boundary
module is free, without division by cycle lengths or group orders.

The global coinvariants vanish by an actual relation. Since
\(v_j=\rho(t_j)v_1\), all their classes agree with \([v_1]\).
The element \(h_0'\) gives
\[
                       (i-1)[v_1]=0.
\]
This forces \([v_1]=0\) in the coinvariants over \(R\).
The dual system has inverse stabilizer character and the same
argument. Full monodromy of the given cover ensures these group
relations are among the relations of its fundamental-group
coefficient complex. In particular both absolute zeroth homology
groups vanish; no central scalar in \(G\) is being assumed.

\paragraph{3. Perfect parabolic pairing and its norm class.}
Apply the bordered-surface cap-product argument of HD-S27 and
HD-S33 to (S34.2). Absolute first homology is free from the graph
complex. Vanishing of absolute zeroth homology makes the
universal-coefficient cross-pairing perfect with relative first
homology, also free. The absolute-to-relative image has free
boundary cokernel and is therefore primitive. The adjoint maps
and extension of functionals from this primitive image prove
surjectivity as well as injectivity of the induced parabolic
duality map. The integral parabolic lattice is consequently
self-dual over \(R\).

Multiplication of its skew-Hermitian cup form by \(i\) gives
the Hermitian form \(H_V\), with the Gaussian convention of
HD-S23. In primitive dimension six it has rank six. Its
determinant is a rational unit of \(R\), and signature
\((3,3)\) makes it \(-2^a\) for some integer \(a\).
Since \(2=N_{K/\mathbb Q}(1+i)\), its ordinary class is
exactly \([-1]\).

\paragraph{4. The actual primitive Jacobian and its Gaussian action.}
Put \(v=(1+i)v_1\), of norm two, and \(D=|G|/m=|H|\).
The matrix-coefficient projector is
\[
 e=\frac m{|G|}\sum_{g\in G}
     \operatorname{Re}h_0(v,\rho(g^{-1})v)\,g.
 \tag{S34.3}
\]
Orthogonality of the matrix coefficients gives
\(e^2=e=e^*\) and \(e\mathbb Q[G]e=K\). These identities
can also be seen directly in the subgroup algebra: (S34.1)
makes the coefficient vanish outside \(H\), and gives
\[
 e=\frac1{|H|}\sum_{h\in H}
           \bigl(\chi(h)^{-1}+\overline{\chi(h)}^{-1}\bigr)h.
 \tag{S34.4}
\]
It is the sum of the two conjugate character idempotents of
\(H\). The induced-module map
\(\operatorname{Ind}_H^G\chi\to V\), sending the coset basis
to the specified \(v_j\), is an actual isomorphism by (S34.1).
Absolute irreducibility therefore makes (S34.4) primitive over
\(\mathbb Q\) with corner \(K\), in agreement with the
matrix-coefficient calculation. All coefficients of \(De\)
are integers, namely \(2\operatorname{Re}\chi(h)\) on \(H\).

Use the actual covariant left deck action \(a_g=L_{g,*}\) on
the full curve Jacobian, and define
\[
 B=\operatorname{im}(De),\qquad
                      i_B=a_{h_0'}^{-1}|_B.
 \tag{S34.5}
\]
This Gaussian action is integral. Indeed (S34.4) is central in
\(\mathbb Q[H]\), so \(h_0'\) commutes with \(e\) and
preserves the actual subvariety \(B\). Moreover
\(e(h_0')^2=-e\), since its two character values are \(i\)
and \(-i\). Thus the square of the restriction is \([-1]\)
on the connected abelian variety \(B\). No central scalar
element of the whole group is needed for (S34.5).
The deck automorphism preserves the restricted polarization,
so its Rosati adjoint is its inverse; the displayed square
therefore gives complex conjugation on this Gaussian action.

The literal coefficient map
\(j_x(g)=h_0(v,\rho(g)x)/2\) and
\(\Phi=\operatorname{PD}(j+\bar j)\) give the paired
isomorphism of HD-S23. The element \(h_0'\) acts on the
coefficient line by \(i\); conjugate linearity in that line
and Poincar\'e duality give
\(\Phi(ix)=a_{h_0'}^{-1}{}_\ast\Phi(x)\).
The complex Schur scalar is \(|G|/(2m)\); the two real cross
terms give precisely
\[
                         h_B(\Phi x,\Phi y)=D H_V(x,y).
 \tag{S34.6}
\]
Its rank-six determinant multiplier is
\(D^6=N_{K/\mathbb Q}(D^3)\). Thus the class computed in
paragraph 3 is the class of the actual restricted-theta factor.

There is also a literal cyclic-cover realization. Let
\(H_0=\ker\chi\), \(X=\widetilde C/H_0\), and
\(Y=\widetilde C/H\). Then \(X\to Y\) is cyclic of degree
four. If \(\alpha:\widetilde C\to X\) is the quotient map,
pullback of its cyclic primitive factor has image \(B\):
in the group algebra
\[
 e=\left(\frac1{|H_0|}\sum_{h\in H_0}h\right)
                              \frac{1-(h_0')^2}{2}.
\]
The pullback and norm identities give this equality on rational
homology, with \(\alpha_*\alpha^*=[|H_0|]\), and pullback
multiplies the principal Jacobian form by \(|H_0|\). Hence it
is an actual polarized isogeny from that cyclic primitive factor
onto \(B\), with the compatible inverse deck Gaussian action.
This proves the relation to the cyclic source rather than merely
assigning the two factors the same dimension.

\paragraph{5. Exact scope and full proof check.}
A different \(K\)-rational coefficient line gives the same
determinant class through a literal nonzero rank-one intertwiner.
If \(e_y\) is its projector, choose \(g\) with
\(h_0(y,\rho(g)v)\ne0\), which exists because the orbit of
\(v\) spans \(V\). Then \(e_yg e\), after clearing its
denominators, is an actual isogeny on the primitive factors;
its line map has positive rational norm ratio and acts as that
line map tensor the identity on the multiplicity space. Its
polarization and determinant changes are the stated rational
similitude and its rank-six norm. For a rank-nine family,
compatible flat-polarization changes under a fixed isogeny also
preserve the class by the explicit common-centralizer result
HD-S17. No arbitrary polarization at an isolated special point
is included in that last assertion.

The full proof above was read by its author. The coset basis,
local cycle quotient, global holonomy relation, primitive
idempotent, integral Gaussian action and Jacobian scalar are
separate exact maps. The result covers arbitrary finite
Gaussian groups with the specified \(K\)-defined line system,
including the actual scalar signed-permutation cases to which
that hypothesis applies. It does not extend a block system
from a larger field or infer a classification of all finite
representations.

The separate \texttt{gaussian\_monomial\_certificate.py}
rebuilds two small weighted-permutation groups directly and
checks the integer subgroup-algebra idempotent, its commutation
with a noncentral \(h_0'\), and its square relation on the
primitive corner. It also checks every weighted local cycle
in those examples. These are finite supplements to the
universal proof above, not classification inputs.
